% ID: SEL-D0031-OA-E04
% ESTADO: cotejada. Segunda pasada visual completa y borrador compilado revisados.
% Fuente: Selectividad/matematicas_aplicadas/titular_junio_6 EXAMEN MAT APLIC C S.pdf
% SHA-256: bebc339cfd8f69c7c4fae160fd165eaef5f6bc93cfd6d15cb9ce289b7029d167
% Página: 1; opción A; ejercicio original 4.

Se desea estimar la media de una variable aleatoria Normal cuya desviación típica es 2.5. Para ello, se toma una muestra aleatoria, obteniéndose los siguientes datos:
\[
18\quad18.5\quad14\quad16.5\quad19\quad20\quad20.5\quad17\quad18.5\quad18
\]
\begin{enumerate}[label=\alph*)]
\item \textbf{(1 punto)} Determine un intervalo de confianza al 96\% para la media poblacional.
\item \textbf{(0.5 puntos)} \textquestiondown{}Cuál es el error máximo cometido con esa estimación?
\item \textbf{(1 punto)} Con el mismo nivel de confianza, si queremos que el error máximo sea inferior a 1, \textquestiondown{}qué tamaño muestral mínimo debemos tomar?
\end{enumerate}
